Table~\ref{tab:s0} summarises S0 over the 30-instance stratified subset. Both violation channels fire on essentially the whole set.

\begin{table}[htbp]
\centering
\caption{S0 violation instrument over 30 instances. Every per-processor route respects the horizon, yet the joint structure does not.}
\label{tab:s0}
\begin{tabular}{lr}
\toprule
Quantity & Value \\
\midrule
Channel A: partially claimed jobs & 184 \;(184 on multi-processor jobs) \\
Mean subgradient norm $\|g\|^2$ at $\mu^\ast$ & 72.5 \\
Channel B: instances cyclic under per-processor optimal order & 21 / 30 \\
Total circuits detected & 148 \\
Acyclic instances with makespan $>L$ & 4 / 9 \\
\quad mean / max excess over $L$ & 1401.0 \;/\; 2044.0 \\
Every per-processor route within $L$ & yes \\
\bottomrule
\end{tabular}
\end{table}

The contrast in the last two rows is the substantive finding. The decomposition certifies every individual route as fitting inside the horizon, while the joint difference-constraint system does not admit a schedule under those same routes. The per-machine subproblem cannot represent the constraint that binds, because that constraint is exactly the one the relaxation dualised.
